\section*{Chapter \ref{ch:iv2:options-and-bond-arithmetic-at-speed} --- Options and Bond Arithmetic at Speed}
\markboth{\protect\omTocMark{}{Answers: Options and Bond Arithmetic at Speed}}{\protect\omTocMark{}{Answers: Options and Bond Arithmetic at Speed}}

\begin{solution}{iq:iv2:options-and-bond-arithmetic-at-speed:1}
About \$8 and about \$4.24. The rule is that an at-the-money call is worth about $0.4\,S\sigma\sqrt{T}$. The exact Black-Scholes values are \$7.97 and \$4.22, so the rule is good to well under 1\% at these volatilities.

Here 0.4 stands for $1/\sqrt{2\pi}$. When rates are zero and the strike is at spot, the Black-Scholes\index{Black-Scholes model} price is $S(2\Phi(\sigma\sqrt T/2) - 1)$, and to first order that is $S \cdot 2\varphi(0) \cdot \sigma\sqrt T/2$. So $0.4 \times 100 \times 0.2 = 8$, and $0.4 \times 50 \times 0.3 \times 0.707 \approx 4.24$.

\iqwrong{Using $\sqrt{T}$ with $T$ in months, or forgetting the square root. Time scales volatility by its square root, so six months is 71\% of a year's volatility, and not half.}
\end{solution}

\begin{solution}{iq:iv2:options-and-bond-arithmetic-at-speed:2}
The straddle costs about \$32, from $0.8\,S\sigma\sqrt T = 0.8 \times 200 \times 0.4 \times 0.5$. The exact value is about \$31.86. It breaks even at about \$32 either way, or 16\%. A typical daily move is about $40\%/16 = 2.5\%$.

The value is twice the call. It is also the expected absolute move, since $E|Z| = \sqrt{2/\pi} \approx 0.8$ standard deviations. So a straddle's break-even, at 0.8 of a standard deviation, is reached a bit more than 40\% of the time.

Divide annual volatility by 16 for a daily move, because $\sqrt{252} \approx 15.9$.

\iqwrong{Dividing by 252 to get a daily move. Variance adds across days, and so volatility scales with the square root of time.}
\end{solution}

\begin{solution}{iq:iv2:options-and-bond-arithmetic-at-speed:3}
The call's delta is about 0.54, and the put's about $-0.46$. The exact values are $\Phi(0.1) \approx 0.5398$ and $-0.4602$.

With zero rates and the strike at spot, $d_1 = \sigma\sqrt T/2 = 0.1$, and $\Phi(x) \approx 0.5 + 0.4x$ for small $x$. So the call's delta is $0.5 + 0.04$. The put's delta is the call's less 1, by put-call parity.

It is above one half because the delta is the chance of finishing in the money under a measure in which the stock is the unit of account, and that gives extra weight to high prices. The median of a lognormal stock is below its mean, and that pulls the ordinary chance below one half: $\Phi(-0.1) \approx 0.46$.

\iqwrong{``An at-the-money option has a delta of exactly 0.5.'' Only when volatility times root time is small, and rates do not lift the forward above the strike.}
\end{solution}

\begin{solution}{iq:iv2:options-and-bond-arithmetic-at-speed:4}
About \$0.40 for the one-year call, about \$0.20 for the three-month call, and about \$2 for the one-year call on the \$500 stock.

The call is worth $0.4\,S\sigma\sqrt T$, so vega\index{vega} is $0.4\,S\sqrt T$ per 1.00 of volatility, and $0.004\,S\sqrt T$ per point. So at \$100 and one year it is 0.40, and at a quarter of a year, half that. It does not depend on the volatility itself, since at the money the price is close to linear in volatility. The exact value at 20\% volatility is about \$0.397.

\iqwrong{Thinking short options have more vega because they are ``more sensitive''. They have more gamma. Vega grows with the square root of time.}
\end{solution}

\begin{solution}{iq:iv2:options-and-bond-arithmetic-at-speed:5}
The put costs about \$7.12. Quickly, $12 - 5 = 7$, plus a small correction to the discount.

By put-call parity\index{put-call parity}, $C - P = S - Ke^{-rT} = 100 - 100e^{-0.05} \approx 4.88$. So $P = 12 - 4.88 = 7.12$.

The quick estimate uses the rate times the strike, 5.00. That overstates by about 0.12, since discounting at 5\% for a year takes off a little less than 5\%.

\iqwrong{Saying that the put and the call at the same strike should cost the same. That is true only when the strike equals the forward, here $100e^{0.05} \approx 105.13$.}
\end{solution}

\begin{solution}{iq:iv2:options-and-bond-arithmetic-at-speed:6}
Duration alone gives a loss of \$487,500. With convexity, the loss is about \$480,750.

The percentage change is $-D\,\Delta y + \tfrac12 C\,\Delta y^2$. So $-6.5 \times 0.30\% = -1.95\%$, which is \$487,500 on \$25 million. The convexity\index{convexity} term is $\tfrac12 \times 60 \times 0.003^2 = 0.027\%$, or \$6,750, and it reduces the loss.

\iqwrong{Subtracting convexity from the price change. Convexity is positive for an ordinary bond, whichever way yields move, and it helps you.}
\end{solution}

\begin{solution}{iq:iv2:options-and-bond-arithmetic-at-speed:7}
The DV01\index{DV01} is \$8,200. You sell about 126 contracts.

The DV01 is value times modified duration times one basis point: $10{,}000{,}000 \times 8.2 \times 0.0001 = 8{,}200$. Then $8{,}200/65 \approx 126.2$.

This hedges a parallel move only. A future's DV01 is that of the bond it tracks, so if that bond sits at a different point on the curve, you still have curve risk.

\iqwrong{Using the bond's face value when the price is far from par. DV01 scales with market value.}
\end{solution}

\begin{solution}{iq:iv2:options-and-bond-arithmetic-at-speed:8}
The price is about 67.6. Macaulay duration is 10, and modified duration is about 9.62. A rise of 10 basis points lowers the price by about 0.65, to about 66.9.

By hand, $e^{-0.4} \approx 0.67$, with $1.04$ close to $e^{0.04}$. The exact value is $100/1.04^{10} = 67.56$. A zero has one cash flow, so its Macaulay duration is its maturity. Modified duration divides by $1 + y$. Then $9.62 \times 0.1\% \times 67.56 \approx 0.65$.

\iqwrong{Taking the 10-year duration as 8 or so, as for a coupon bond. There are no coupons to pull the average life in.}
\end{solution}

\begin{solution}{iq:iv2:options-and-bond-arithmetic-at-speed:9}
Quick answer 5\%. The exact answer is about 5.01\%.

Investing for two years at 4\% must match investing for one at 3\% and then one at the forward rate: $1.04^2 = 1.03(1 + f)$, so $f = 1.0816/1.03 - 1 \approx 5.01\%$.

The quick answer works because logarithms add: the two-year rate is about the average of the one-year rate and the forward, so $f \approx 2 \times 4\% - 3\%$.

\iqwrong{Averaging 3\% and 4\%. The two-year rate is itself the average, and the forward is what makes it so.}
\end{solution}

\begin{solution}{iq:iv2:options-and-bond-arithmetic-at-speed:10}
Quickly, $1.10 \times 1.02 = 1.122$. Exactly, $1.10 \times 1.05/1.03 \approx 1.1214$. At 1.13 the forward is too high, so sell euros forward and buy them today with borrowed dollars.

The higher-rate currency trades at a forward discount, which here means more dollars per euro. Do the lock per euro delivered in a year: borrow dollars today, buy $1/1.03$ euros at 1.10 for about \$1.068, and put them on deposit. In a year, you hold 1 euro, which you deliver at 1.13. You owe $1.068 \times 1.05 \approx \$1.1214$. Your profit is about 0.86 cents per euro.

This is covered interest parity\index{covered interest parity}.

\iqwrong{Thinking that the higher-rate currency should strengthen in the forward. The forward offsets the rate gap, so the dollar with the higher rate trades weaker forward against the euro.}
\end{solution}

\begin{solution}{iq:iv2:options-and-bond-arithmetic-at-speed:11}
Buy the basket, create 50,000 ETF shares and sell them at the bid. You make \$0.09 per share, or \$4,500 per unit.

Selling at \$50.02 against a total cost of $\$49.90 + \$0.03 = \$49.93$ leaves \$0.09. The ETF is rich, so you sell it and create it. Going the other way, buying the ETF at \$50.04 and redeeming for a basket worth \$49.90, loses money.

\iqwrong{Comparing the basket with the ETF's mid or ask. You sell at the bid, and costs come off first.}
\iqfollow{What is at risk?}{The prices can move while you trade, and the basket may not be fully tradable at the quoted value. So you buy the basket and sell the ETF at the same moment, and create later.}
\end{solution}

\begin{solution}{iq:iv2:options-and-bond-arithmetic-at-speed:12}
Implied volatility is about 30\%. The one-standard-deviation move to expiry is about \$15.

The straddle costs $0.8\,S\sigma\sqrt T$, so $\sigma = 12/(0.8 \times 100 \times 0.5) = 30\%$. One standard deviation over three months is $100 \times 0.3 \times 0.5 = 15$, and $12/0.8 = 15$ gives the same.

So a straddle's price is about 0.8 of the standard deviation, and a standard deviation is about 1.25 straddle prices.

\iqwrong{Taking the straddle's price as the standard deviation. It is the expected absolute move, which is smaller.}
\end{solution}

\begin{solution}{iq:iv2:options-and-bond-arithmetic-at-speed:13}
Gamma is about 0.02 per dollar. Theta is about 1.6 cents a day. To cover it the stock must move about \$1.26, or about 1.26\%, which is the daily move at 20\% volatility.

Gamma\index{gamma} is $\varphi(d_1)/(S\sigma\sqrt T) \approx 0.4/(100 \times 0.2) = 0.02$. The exact value is about 0.0198.

The first route to theta\index{theta} is Black-Scholes with zero rates: $\theta = -\tfrac12\Gamma S^2\sigma^2$, which is $0.5 \times 0.02 \times 10{,}000 \times 0.04 = 4$ a year, or $4/252 \approx 0.016$ a day. The second route takes the rule $C \approx 0.4S\sigma\sqrt T$: the derivative in $T$ gives $C/(2T)$, so $8/2/252 \approx 0.016$.

Gamma's gain on a move $m$ is $\tfrac12\Gamma m^2$, and setting this equal to theta gives $m^2 = S^2\sigma^2/252$. So the break-even move equals one day's standard deviation, $20/\sqrt{252} \approx 1.26$.

\iqwrong{Thinking that theta and gamma are separate facts. For a delta-hedged option, they are two sides of the same bet on the size of the daily move.}
\end{solution}

\begin{solution}{iq:iv2:options-and-bond-arithmetic-at-speed:14}
The call less the put is about \$4.90. Without the dividend, it would be about \$6.88.

With a dividend, $C - P = S - De^{-rt} - Ke^{-rT}$. The dividend is discounted for a quarter of a year at 4\%, a factor of about 0.99, so about 1.98. The strike is discounted for half a year, a factor of about 0.98, so about 93.12. Then $100 - 1.98 - 93.12 = 4.90$. Without the dividend, it is $100 - 93.12 = 6.88$.

\iqwrong{Taking the dividend off at its full amount and not discounting it. That is small here, but the bigger error is forgetting it when the options span the date: parity is then off by almost \$2.}
\end{solution}

\begin{solution}{iq:iv2:options-and-bond-arithmetic-at-speed:15}
The options imply a forward of about \$100.51. Carry alone implies \$101.50. So there is a gap of \$1, or about 2\% a year, which could be a dividend, a stock-borrow cost, or an error in the rate.

Long a call and short a put at strike $K$ is a forward at $K$, so the forward less $K$, discounted, is $C - P$. With simple rates, $0.50 \times 1.015 \approx 0.51$, so the forward is about \$100.51. Carry alone gives $100 \times (1 + 3\% \times 0.5) = 101.50$.

The stock is trading at a lower forward than carry says. So someone pays to hold the stock forward less than it costs to carry it. If you can borrow the stock cheaply and there is no dividend, you could buy the synthetic and sell the stock.

\iqwrong{Calling it an arbitrage immediately. A gap of this size is usually a dividend you missed or a hard-to-borrow stock, and in a hard-to-borrow stock you cannot do the arbitrage cheaply.}
\end{solution}

\begin{solution}{iq:iv2:options-and-bond-arithmetic-at-speed:16}
Its price is 100 because the coupon equals the yield. Its modified duration equals the annuity factor $(1 - 1.05^{-10})/0.05 \approx 7.72$. So the price falls about 0.77 to about 99.23, and the exact price is 99.23.

A coupon paid at the same rate you discount at is exactly what the principal earns, so the bond is worth par. Differentiate the price with respect to the yield, and at par, the result is the present value of a 1 annuity. So a par bond's modified duration is the annuity factor.

\iqwrong{Using 10 as the duration. A coupon bond's cash arrives before maturity, so its duration is shorter, and it falls as the coupon rises.}
\end{solution}

\begin{solution}{iq:iv2:options-and-bond-arithmetic-at-speed:17}
Carry is 0.25 points, and roll-down is about 0.11 points, for a total of about 0.36. The yield can rise about 8 basis points before you lose money.

Carry\index{carry} is the yield less the financing rate, times the time: $(4.5\% - 3.5\%) \times 0.25 = 0.25$. Roll-down\index{roll-down} comes from ageing. After three months, the bond has 4.75 years left, and on an unchanged curve it yields 2.5 basis points less. That gains $4.5 \times 0.025 = 0.1125$.

The break-even rise is $0.3625/4.5 \approx 8$ basis points.

\iqwrong{Counting the coupon as carry without taking off the financing. Financing is what makes a position with no capital cost anything.}
\end{solution}

\begin{solution}{iq:iv2:options-and-bond-arithmetic-at-speed:18}
The published value shows a premium of 1.25\%. The real one is about 0.25\%, on a fair value of \$40.40. Your cost is 20 basis points, so about 5 basis points remains, and it is at best a marginal trade.

Tokyo is shut, so the published net asset value\index{net asset value} is stale. Futures have moved 1\%, so the better estimate of the basket is $40 \times 1.01 = 40.40$. Then $40.50/40.40 - 1 \approx 0.25\%$.

\iqwrong{Selling the ETF because it ``trades 1.25\% above its value''. The ETF is reflecting news that the published number cannot, and it can do so only because it trades while Tokyo is closed.}
\end{solution}

\begin{solution}{iq:iv2:options-and-bond-arithmetic-at-speed:19}
Fair value is 4,030. The future is 6 points rich, or 15 basis points of spot. Sell the future and buy the basket. After costs of 10 basis points, 5 remain, or 2 points.

Fair value is spot plus financing less dividends: $4000 \times (1 + (5\% - 2\%) \times \tfrac14) = 4030$. The 6-point gap is $6/4000 = 15$ basis points. You hold this cash and carry\index{cash and carry} to expiry, and the futures and basket then converge.

\iqwrong{Quoting the gap without its costs and its risks. Dividends are estimates, and a shift in them or in your financing can be larger than the 5 basis points left.}
\end{solution}

\begin{solution}{iq:iv2:options-and-bond-arithmetic-at-speed:20}
The price says a chance of about 37\% of finishing above about \$105. The Black-Scholes value is about 36.5\%.

A call spread pays 0 below 100 and the full 10 above 110, and it is linear in between. So its value per unit of width is roughly the risk-neutral chance of finishing beyond the middle, a digital\index{digital option} at 105. And $3.67/10 = 0.367$. For the model's value, $d_2 = \ln(100/105)/0.2 - 0.1 = -0.344$, and $\Phi(-0.344) \approx 36.5\%$.

So you read a chance off two call prices.

\iqwrong{Calling that the real-world chance. It is risk-neutral, and it includes risk premia and the volatility smile.}
\end{solution}

\begin{solution}{iq:iv2:options-and-bond-arithmetic-at-speed:21}
The normal estimate is about \$3.96, and Black-Scholes gives about \$4.29. It is lower because a lognormal stock has a fatter right tail than a normal with the same standard deviation.

The normal formula, the Bachelier\index{Bachelier model} model, is $C = s\,(\varphi(d) - d(1 - \Phi(d)))$, with $d = (K - S)/s = 0.5$ and $s = 20$. Then $20 \times (0.35 - 0.5 \times 0.31) \approx 3.9$, and with four figures it is 3.96.

The exact Black-Scholes value is 4.29. The difference of about 8\% shows how far a one-year 20\% option is from the normal limit.

\iqwrong{Drawing a straight line from the at-the-money price using its delta of about one half: $8 - 5 = 3$. A call is convex in the strike, so its value falls more and more slowly as you leave the money, and the straight line comes out too low.}
\end{solution}

\begin{solution}{iq:iv2:options-and-bond-arithmetic-at-speed:22}
The rule of thumb says about 0.16, since $\Phi(-1) \approx 0.16$. The exact delta is about 0.21 with one year and 20\%, and about 0.17 for the short option with 5\%. They differ because a lognormal distribution has its centre and its width on a different scale from a normal one, and that matters more as total volatility grows.

Delta is $\Phi(d_1)$, with $d_1 = \ln(S/K)/(\sigma\sqrt T) + \tfrac12\sigma\sqrt T$. At 120 and 20\%, $d_1 = \ln(100/120)/0.2 + 0.1 \approx -0.81$, which gives 0.21. At 105 and 5\%, $d_1 = \ln(100/105)/0.05 + 0.025 \approx -0.95$, which gives 0.17.

The logarithm of 1.2 is only 0.182, and $\tfrac12\sigma\sqrt T$ lifts it further. So at higher total volatility, a strike that is 1 standard deviation away in dollars is closer than that in logs. The rule fits low total volatility.

\iqwrong{Using the rule to set a strike for a precise hedge. It is a rough guide for short dates.}
\end{solution}

\begin{solution}{iq:iv2:options-and-bond-arithmetic-at-speed:23}
Sell about \$12.1 million of the 2-year and about \$2.77 million of the 10-year. The butterfly's spread is 20 basis points. If only the 5-year falls 5 basis points, you make \$23,000. If all three rise 10, you make nothing, to first order.

The 5-year position has a DV01 of $10 \times 460 = \$4{,}600$. Each wing gets half: $2{,}300/190 \approx 12.1$ million, and $2{,}300/830 \approx 2.77$ million. The net is zero, so parallel moves cancel. Half on each wing also means that a pure steepening cancels, where the wings move in opposite directions.

The spread is $2 \times 4.4 - 4.0 - 4.6 = 0.2$. On a fall in the 5-year alone, you make $4{,}600 \times 5$.

\iqwrong{Matching notional instead of risk. Selling \$5 million of each wing leaves $4{,}600 - 950 - 4{,}150 = -500$, a net short of \$500 of DV01, which is a bet on the curve's level.}
\iqfollow{Is this trade cash-neutral?}{No. You buy \$10 million and sell about \$14.9 million of face value, and so you have to finance the difference.}
\end{solution}

\begin{solution}{iq:iv2:options-and-bond-arithmetic-at-speed:24}
You collect about \$6.91. A move of 1\% a day is about 16\% volatility. Your expected profit is about \$3.2. That is the straddle at 30\% less the straddle at 16\%, and not a figure from a constant gamma.

The straddle is $0.8 \times 100 \times 0.3 \times \sqrt{1/12} \approx 6.93$, and the exact value is 6.91. At 16\%, the straddle is about 3.68. Their difference is about 3.2.

A rule says vega times the volatility gap. Straddle vega is $0.8 \times 100 \times 0.2887 = 23.1$ per 1.00, and times 0.14 is 3.23. This holds because at the money, the price is close to linear in volatility.

Now take gamma held constant. The straddle's gamma is about 0.092, and $\tfrac12\Gamma S^2(\sigma_i^2 - \sigma_r^2)T$ gives only about 2.47. This understates because gamma is not constant: in the last days, it rises for a stock that stays near the strike. The exact statement is $V(\sigma_i) - V(\sigma_r) = E\left[\int \tfrac12\Gamma_i S^2(\sigma_i^2 - \sigma_r^2)\,dt\right]$, by Itô's lemma\index{Ito's lemma} on the option's value.

A single path can be far from that mean, because the gain depends on where the stock has been when the moves happen.

\iqwrong{Claiming a sure profit of 3.2. Hedging at implied volatility gives a mean, not a certain gain. Only hedging at the true realised volatility makes the profit certain, and you do not know that when you start.}
\end{solution}

\begin{solution}{iq:iv2:options-and-bond-arithmetic-at-speed:25}
The straddle implies a total standard deviation of 10\%. Five days at 30\% have a variance of about $0.0018$, which leaves an earnings move of about 9.1\%. If you thought the move would be about 6\%, you would sell the straddle and hedge.

The total standard deviation is $8/0.8 = 10$ points, or a variance of 0.01. Ordinary days have a variance of $0.3^2/252 \approx 0.000357$ each, or 0.00179 for five. The rest is $0.00821$, and its root is about 9.06\%.

Variance\index{event volatility} adds. Standard deviations do not, so subtracting 5 days' 4.2\% from 10\% would give 5.8\%, which is wrong.

At your 6\%, total variance is $0.0036 + 0.00179 = 0.00539$, a standard deviation of 7.34\%. A fair straddle is then about $0.8 \times 7.34 \approx \$5.87$, against the \$8 you can sell at.

\iqwrong{Dividing the straddle price by the number of days, or forgetting to take the ordinary days out. The earnings day carries about 82\% of the variance here.}
\iqfollow{What are the risks?}{The historical earnings moves may be lower than the next, and the stock can gap through your hedges. Size so that a gap of three times your 6\% is survivable.}
\end{solution}
